Seja a eq. (3.6) uma lei de potência da forma
\[
y = C\,x^{n},
\]
onde $C > 0$ é a constante de proporcionalidade e $n$ é o expoente de escala.

Aplicando o logaritmo (base 10) a ambos os membros:
\[
\log_{10} y = \log_{10} C + n\,\log_{10} x.
\]
Definindo as variáveis transformadas $Y = \log_{10} y$ e $X = \log_{10} x$, obtém-se a equação de uma reta:
\[
Y = n\,X + \log_{10} C.
\]

Essa é a \textbf{forma linear para gráfico log-log} da eq. (3.6). No plano log-log:
\begin{itemize}
\item a \textbf{inclinação} da reta é o expoente $n$;
\item o \textbf{intercepto vertical} é $\log_{10} C$, do qual se extrai a constante $C = 10^{\mathrm{intercepto}}$.
\end{itemize}

Portanto, um gráfico log-log transforma a lei de potência em uma reta cujos parâmetros identificam diretamente o expoente e a amplitude da equação original.

